% docs/texte/10_engineering/api_docs.tex -- chapter "API Documentation" (Interferosim)
% Version: 0.1.1 (2026-10-12) -- script build_api_docs
%   0.1.0: first version
% ================================================================

\index{Doxygen}\index{API documentation}
The documentation of the source code (modules, classes, functions) is
generated with Doxygen from the docstrings of the package
\texttt{src/interferosim}. The configuration is the file
\texttt{Doxyfile} in the project root.

\section{Installation}

\begin{itemize}
  \item \textbf{Doxygen} (required): Windows installer from
        \url{https://www.doxygen.nl/download.html}; the directory
        \texttt{\ldots\textbackslash doxygen\textbackslash bin} must be on
        the \texttt{PATH} (the installer offers this). Linux:
        \texttt{sudo apt install doxygen}.
  \item \textbf{Graphviz} (optional): for collaboration and call graphs
        (\url{https://graphviz.org}); then set \texttt{HAVE\_DOT = YES}
        in the \texttt{Doxyfile}. Without Graphviz Doxygen draws the
        class diagrams itself.
  \item \textbf{\LaTeX} (optional): for the PDF reference manual
        (MiKTeX is sufficient).
\end{itemize}

\section{Generating the documentation}

\paragraph*{With the script (Windows).} \texttt{build\_api\_docs.bat}
(double-click) or \texttt{build\_api\_docs.ps1} finds Doxygen, installs
it with \texttt{winget} if it is missing (package
\texttt{DimitriVanHeesch.Doxygen}), generates the documentation and opens
the start page. Options: \texttt{-Pdf} (also the PDF reference manual),
\texttt{-NoOpen}, \texttt{-NoInstall}.

\paragraph*{By hand.} In the project folder:
\begin{lstlisting}[basicstyle=\ttfamily\footnotesize, frame=single]
doxygen Doxyfile                       # HTML and LaTeX into docs/api/
start docs\api\html\index.html         # open in the browser (Windows)

cd docs\api\latex                      # optional: PDF reference manual
make.bat                               # -> refman.pdf
\end{lstlisting}

\begin{center}
\begin{tabular}{@{}ll@{}}
\toprule
Output & Content \\
\midrule
\texttt{docs/api/html/index.html} & HTML with tree view, search, class
  diagrams, source browser \\
\texttt{docs/api/latex/refman.pdf} & PDF reference manual (after
  \texttt{make.bat}) \\
\texttt{docs/api/doxygen\_warnings.txt} & warnings of the last run \\
\bottomrule
\end{tabular}
\end{center}

The folder \texttt{docs/api/} is generated and therefore excluded from
version control (\texttt{.gitignore}). The main page is the project
\texttt{README.md}. The version shown (\texttt{PROJECT\_NUMBER}) is
raised together with the package version.

\section{Docstrings}

Doxygen takes the Python docstrings as documentation
(\texttt{PYTHON\_DOCSTRING = YES}):
\begin{itemize}
  \item every module starts with a docstring describing its content,
        origin and version history;
  \item every public class and function has a docstring; the first
        sentence is the short description shown in the overview lists
        (\texttt{JAVADOC\_AUTOBRIEF = YES});
  \item parameters, units and conventions are described in the
        docstring (e.\,g.\ \enquote{u, v in wavelengths}, \enquote{Jy per
        pixel}).
\end{itemize}
